% =============================================================================
% Bundled Fonts and Default-Typeface Changes in Major Operating Systems,
% Office Suites, and Document Tooling: A Provenance Timeline Reference
% XeLaTeX + xeCJK, two-column  (Overleaf: set the compiler to XeLaTeX)
% Build: xelatex -> bibtex -> xelatex x2   (see Makefile)
% CJK fonts required (macOS system fonts): Songti TC (main), Hiragino Mincho ProN
% (Japanese, via \ja), Apple SD Gothic Neo (Korean, via \ko).
% =============================================================================
\documentclass[10pt,a4paper]{article}

% --- Packages ----------------------------------------------------------------
\usepackage{fontspec}
\usepackage{xeCJK}
\setCJKmainfont{Songti TC}
\setCJKsansfont{Songti TC}
\setCJKmonofont{Songti TC}
\newCJKfontfamily\japanesefont{Hiragino Mincho ProN}
\newCJKfontfamily\koreanfont{Apple SD Gothic Neo}
\usepackage{amsmath,amssymb,amsfonts}
\usepackage{booktabs}
\usepackage{array}
\usepackage{multirow}
\usepackage{graphicx}
\usepackage{xcolor}
\usepackage{hyperref}
\usepackage[margin=0.75in]{geometry}
\usepackage{enumitem}
\usepackage{float}
\usepackage{caption}
\usepackage{longtable}
\usepackage[numbers,sort&compress]{natbib}
\usepackage{xurl}
\usepackage{authblk}
\usepackage{tikz}
\usetikzlibrary{positioning,calc,arrows.meta,fit,backgrounds}
\usepackage{textcomp}
\usepackage{adjustbox}
\usepackage[framemethod=tikz]{mdframed}

\setlength{\columnsep}{18pt}
\setlength{\emergencystretch}{3em}

\hypersetup{
    colorlinks=true,
    linkcolor=blue!70!black,
    citecolor=green!50!black,
    urlcolor=blue!60!black,
}

% --- Script helpers ----------------------------------------------------------
% Songti TC (main CJK) covers Simplified and Traditional Chinese. Japanese kana
% and Korean hangul are not in Songti TC, so per-script macros switch fonts.
\newcommand{\ja}[1]{{\japanesefont #1}}
\newcommand{\ko}[1]{{\koreanfont #1}}
\newcommand{\code}[1]{\texttt{#1}}
% Japanese title held in a macro so it renders in the Japanese font from the .bbl.
\newcommand{\ichitarotitle}{\ja{「一太郎2021」2月5日発売}}

% --- Aeon-Timeline-style lane construct --------------------------------------
% A vertical spine (tikz rule via the array `!{}' separator) with year nodes on
% the left and event labels on the right. Robust in one- or two-column context;
% pass an optional label-column width for figure* (full-width) timelines.
\definecolor{fpspine}{HTML}{1B9E8A}
\newcommand{\fpyear}[1]{%
  \tikz[baseline=(fpy.base)]{\node[fill=fpspine,text=white,rounded corners=2pt,
    inner sep=2.5pt,font=\footnotesize\bfseries] (fpy) {#1};}}
\newenvironment{fptimeline}[1][0.80\textwidth]{%
  \renewcommand{\arraystretch}{1.3}%
  \begin{tabular}{@{}r@{\hspace{6pt}}!{{\color{fpspine}\vrule width 1.3pt}}@{\hspace{9pt}}>{\raggedright\arraybackslash}p{#1}@{}}}%
  {\end{tabular}}
\newcommand{\fpevent}[2]{\fpyear{#1} & #2\\}

% --- Title -------------------------------------------------------------------
\title{%
  \textbf{Bundled Fonts and Default-Typeface Changes in Major Operating Systems,
  Office Suites, and Document Tooling:}\\[4pt]
  \large A Provenance Timeline Reference%
}
\author{Albert Hui}
\date{July 24, 2026 \\[2pt] \small (working paper)}

% =============================================================================
\begin{document}

\twocolumn[
\maketitle

\begin{abstract}
The set of fonts bundled with an operating system or application, and the
\emph{default} typeface each release binds to its user interface and to new
documents, change at datable public releases. Those dates bound the earliest time
an artifact could have acquired the font through ordinary public channels
(pre-release access---foundries, beta programs, OEMs---is the caveat), and
default-font choices identify candidate producing software. This paper assembles
verified timelines of (1)~default system/UI font changes and bundled-font
milestones in Windows, macOS, iOS, Android, and desktop Linux, including per-script
CJK lanes (Simplified Chinese, Traditional Chinese, Hong Kong, Japanese, Korean);
(2)~default document-font changes in major office suites (Microsoft Office
including its East Asian editions, OpenOffice.org/LibreOffice, WPS Office, Google
Docs, Apple iWork, JustSystems Ichitaro, Hancom Office); (3)~default and
substitution fonts in PDF and scan/OCR tooling (Adobe Acrobat/Reader, PaperPort and
OCR peers); and (4)~default-font and font-substitution behavior in
document-generation libraries (the Aspose suites and peers). Events are cited to the
most authoritative sources located, with hotlinked URLs, and citations are annotated
with verification tiers; where only secondary or incomplete sourcing was found, the
citation and the claim say so, and claims that could not be adequately verified are
labeled UNCERTAIN.
\end{abstract}
\vspace{12pt}
]

% === §1 Introduction =========================================================
\section{Introduction}\label{sec:intro}

A font cannot appear in a document before the font is available. The canonical
demonstration is the Calibri affair of 2017 (``Fontgate''), in which property
documents represented as executed in \textbf{February 2006} were typeset in
Calibri---a font that first became publicly obtainable with Office 2007 Beta 2 on
2006-05-23 and generally available with Windows Vista and Office 2007 on
2007-01-30. The February 2006 dating is the decisive detail: the documents predate
even the public beta, which featured in Pakistan's Supreme Court proceedings against
the Sharif family.\cite{intro-fontgate} Beyond first-availability bounds,
\emph{default} fonts carry a second, weaker signal: they identify the factory
configuration of candidate producing software. This reasoning is asymmetric. A font
used before its earliest public availability is decisive (an anachronism); a
default-font match is only \emph{consistent with} a toolchain, never proof of it,
because users change defaults, templates override them, and fonts travel. Read this
way: a PDF whose text layer is set in Fanwood, accompanied by Aspose's
characteristic substitution warning or metadata, is strong evidence of Aspose.Words
rendering on a font-starved Linux host (\S\ref{sec:libraries}); text in Adobe Sans
MM is consistent with Acrobat rendering a PDF whose fonts were missing
(\S\ref{sec:pdf}); a Simplified-Chinese document body in 等线 (DengXian) cannot
predate 2015-07-29 (the font's first ship, in Windows 10) and is consistent with the
Office 2016-era Simplified-Chinese default path (\S\ref{sec:msoffice}).

This paper is a dated, sourced reference for such reasoning. Each platform or product
gets an Aeon-Timeline-style lane view plus an event register with citations.

% === §2 Method and evidence tiers ============================================
\section{Method and Evidence Tiers}\label{sec:method}

\paragraph{Search protocol.} Research was conducted in July 2026 by parallel
per-platform investigations, each proceeding: (1)~general web search to locate
candidate sources for each event; (2)~preference-ordered source selection---vendor
primary documentation (Microsoft Learn/Typography, Apple Newsroom/Support/Developer,
Android Developers Blog, project release notes and bug trackers, vendor API
references), then institutional records (Computer History Museum, SEC EDGAR), then
first-person designer accounts (Folklore.org; Ken Lunde's CJKV writings treated as
near-primary), then reputable secondary/encyclopedic sources; (3)~fetch-and-confirm
verification of each cited URL against the specific claim (a 200 response serving
unrelated content does not count as verification); (4)~Wayback Machine fallback for
dead originals. Searches used general-purpose web search engines plus direct
navigation of the vendor documentation trees named above; fetching used both plain
HTTP clients and rendering fetchers (several vendor sites are JavaScript-routed or
bot-hostile, noted per citation). Inclusion required that a source bear directly on
the dated event; user-forum material was admitted only for vendor-staff answers or
where no better source exists, and is always tiered [secondary]. Artifact
inspections (e.g., \S\ref{sec:msoffice}'s theme dump) are documented with path,
software version, observation date, and SHA-256. Conflicting sources are footnoted
with both readings and a stated preference (vendor primary over encyclopedic---e.g.,
the Trebuchet MS first-ship discrepancy in \S\ref{sec:win-latin}). Negative findings
(``no change found'') are scoped to the sources actually searched and say so. The
research window bounds all ``current'' statements. This protocol supports
spot-checking every claim from its citation; it is not a systematic-review protocol,
and \S\ref{sec:limitations} says so.

\paragraph{What counts as a font ``appearing''.} Provenance reasoning must
distinguish: (a)~a \emph{nominal reference}---a font name in document markup (a
\code{.docx} run property, a PDF \code{/BaseFont} name) with no font data; (b)~an
\emph{embedded font program}; (c)~\emph{rendered output}---glyph outlines actually
drawn, possibly by a substituted face; (d)~a \emph{substituted display face} (the
viewer's stand-in, \S\ref{sec:pdf}); and (e)~an \emph{OCR text layer's} styling,
which reflects the OCR engine, not the source document
(\S\ref{sec:ocr}). Anachronism logic applies strictly to (b)---an embedded font
program cannot predate the font. For (a), a nominal reference is strong but
rebuttable: font \emph{names} can be written by software that never possessed the
font (template inheritance, format converters, substitution tables---e.g., 2007-era
iPhone pages could \emph{request} Helvetica Neue years before the font shipped on the
device, \S\ref{sec:ios-latin}), so a nominal anachronism obliges the examiner to
rule out name-injection paths before concluding. Inferences from (c)--(e) concern the
\emph{rendering/processing} toolchain, not the original author. The timelines below
date public availability; each anchor's type is stated in \S\ref{sec:anchors}.

\paragraph{Evidence tiers.} Citations below carry one of these statuses, assigned
during research in July 2026:
\begin{itemize}[leftmargin=1.3em,itemsep=1pt]
  \item \textbf{[verified live]}---the URL was fetched during this research and its
        content was confirmed to support the claim (an HTTP 200 serving an error or
        unrelated page does not qualify).
  \item \textbf{[archive]}---the original is dead or unreachable; an Internet Archive
        snapshot is cited. Where the snapshot's \emph{existence} was confirmed via
        the Wayback availability API but its body could not be re-read during this
        research (archive.org rate-limiting), the citation says so explicitly.
  \item \textbf{[secondary]}---a reputable non-primary source (encyclopedic or
        press), used for corroboration or where no primary survives; presented as
        such.
  \item \textbf{UNCERTAIN}---the claim is plausible and commonly reported but could
        not be verified against an adequate source during this research; it is
        flagged, not asserted.
\end{itemize}

Primary sources are preferred throughout: vendor documentation (Microsoft
Learn/Typography font-family pages, Apple Support ``Fonts included with\ldots''
articles, Apple Newsroom press releases, Android Developers Blog,
GNOME/KDE/Fedora/Debian/Ubuntu project records, Adobe HelpX and the PDF/ISO
specification, vendor API references), institutional records (Computer History
Museum), and first-person designer accounts (Folklore.org; Ken Lunde's CJK writings
are treated as near-primary for CJKV matters).

\paragraph{Script lanes, not retail regional editions.} For CJK coverage this paper
models \emph{script/locale lanes}---Simplified Chinese (SC), Traditional Chinese
(TC), Traditional Chinese--Hong Kong (TC-HK), Japanese (JP), Korean (KR)---rather
than retail regional editions (a ``Taiwan edition'' vs a ``mainland edition''). Three
facts justify this choice. First, since the Unicode consolidation era (Windows 2000,
Mac OS X, and all mobile OSes from birth), every edition of each OS carries the
\emph{same worldwide font set}; what varies by locale is the per-locale default UI
font binding and, in office software, the per-language document
defaults.\cite{m-w2k,m-w2k-2} Second, regional editions matter as distinct artifacts
only in the pre-Unicode era (Chinese/Japanese/Korean Windows 3.x/9x; KanjiTalk-era
classic Mac OS), and those appear here as dated events \emph{within} the script
lanes. Third, Singapore has no distinct font lane---zh-SG follows Simplified Chinese
conventions---while Hong Kong earns a sub-lane solely because of the HKSCS
character-set supplements and, later, dedicated HK typefaces (PingFang HK,
\S\ref{sec:apple-cjk}).\cite{m-hkscs}

A note on Windows 2000 specifically, because it anchors the ``editions stop
mattering'' argument: on Windows 2000 and XP the East Asian fonts shipped on the
installation media of every edition but were \textbf{not installed by default}
outside CJK locales (Regional Options $\rightarrow$ ``Install files for East Asian
languages''); CJK fonts became installed-by-default regardless of locale only in the
Windows Vista/7 era.\cite{m-w2k-install}

% === §3 Operating systems ====================================================
\section{Operating Systems}\label{sec:os}

\subsection{Microsoft Windows --- Latin lane}\label{sec:win-latin}

Figure~\ref{fig:win-latin} traces the Windows default UI font and bundled-font
milestones.

\begin{figure*}[t]
\centering
\begin{fptimeline}
\fpevent{1985}{Windows 1.0 (Nov 20)---bitmap System font (titles and menus) with Helv bitmap for dialogs}
\fpevent{1992}{Windows 3.1 (Apr 6)---TrueType arrives with Arial, Times New Roman, Courier New; Helv renamed MS Sans Serif}
\fpevent{1995}{Windows 95 (Aug 24)---MS Sans Serif as default UI face; Comic Sans MS via Plus! pack}
\fpevent{1998}{Windows 98 (Jun 25)---Verdana, Tahoma, Trebuchet MS present per Microsoft font tables}
\fpevent{2000}{Windows 2000 (Feb 17)---Tahoma replaces MS Sans Serif as default UI font; Georgia bundled}
\fpevent{2001}{Windows XP (Oct 25)---Tahoma retained}
\fpevent{2007}{Windows Vista (Jan 30)---Segoe UI 9\,pt becomes default; ClearType six (Calibri, Cambria, Candara, Consolas, Constantia, Corbel) ship}
\fpevent{2012}{Windows 8 (Oct 26)---Segoe UI stylistic refresh, Semilight weight}
\fpevent{2015}{Windows 10 (Jul 29)---Segoe UI retained; many script fonts move to optional Feature-on-Demand packs}
\fpevent{2021}{Windows 11 (Oct 5)---Segoe UI Variable becomes default; Cascadia Code and Cascadia Mono bundled}
\fpevent{2023}{Aptos (Jul 13)---replaces Calibri as Microsoft 365 default; distributed via Office cloud fonts, never bundled in Windows}
\end{fptimeline}
\caption{Windows default UI font and bundled-font milestones (Latin lane).}
\label{fig:win-latin}
\end{figure*}

\paragraph{Event register.}
\begin{itemize}[leftmargin=1.2em,itemsep=2pt]
\item \textbf{1985-11-20 --- Windows 1.0.} UI set in bitmap fonts: a proportional
``System'' font for titles/menus and ``Helv'' for dialogs; no scalable
outlines.\cite{w-mssans}
\item \textbf{1992-04-06 --- Windows 3.1.} TrueType introduced, with the founding
scalable core set \textbf{Arial, Times New Roman, Courier New} (plus Symbol)---the
first Windows with scalable outline fonts. The ``Helv'' bitmap is renamed \textbf{MS
Sans Serif}. Launch date per the Computer History
Museum.\cite{w-31date,w-ttf,w-ttf-2,w-mssans}
\item \textbf{1995-08-24 --- Windows 95.} \textbf{MS Sans Serif} continues as the
default UI font (through NT 4.0, 98, ME). \textbf{Comic Sans MS} first ships with the
Windows 95 Plus! pack and OEM releases, per Microsoft's own font
page.\cite{w-95date,w-mssans,w-comic}
\item \textbf{1998-06-25 --- Windows 98.} Microsoft's per-font supply tables list
\textbf{Verdana} (v2.10), \textbf{Tahoma} (v2.25), and \textbf{Trebuchet MS} (v0.98)
from Windows 98. Verdana and Georgia (both Matthew Carter, 1996) reached users
earlier via Internet Explorer and the 1996 ``Core fonts for the Web'' program;
Wikipedia's typeface list attributes Verdana/Tahoma to Windows 95 via that route,
while Microsoft's own tables enumerate only from 98---and for Trebuchet MS
Microsoft's page (Windows 98) contradicts Wikipedia's ``2000''. Microsoft's pages are
preferred here; the discrepancies are noted rather than
resolved.\cite{w-98date,w-verdana,w-tahoma,w-trebuchet,w-corefonts,w-typelist}
\item \textbf{2000-02-17 --- Windows 2000.} \textbf{Tahoma replaces MS Sans Serif} as
the default screen/UI font (through XP and Server 2003). \textbf{Georgia} appears
from Windows 2000 in Microsoft's table.\cite{w-2000date,w-tahoma-default,w-georgia}
\item \textbf{2001-10-25 --- Windows XP.} Tahoma retained as default UI
font.\cite{w-xpdate,w-tahoma-default}
\item \textbf{2007-01-30 --- Windows Vista (consumer GA; business availability
2006-11-30).} \textbf{Segoe UI at 9\,pt becomes the default UI font}, per Microsoft's
Win32 UX guidance. The \textbf{ClearType Font Collection}---Calibri, Cambria, Candara,
Consolas, Constantia, Corbel (plus Cambria Math and Meiryo)---ships to the general
public; it had first appeared publicly in Office 2007 Beta 2 on 2006-05-23. In Office
2007, released the same day, \textbf{Calibri replaced Times New Roman as the Word
default} and Arial as the Excel/PowerPoint/Outlook default (\S\ref{sec:msoffice}).
Provenance anchor: a document purportedly authored before mid-2006 set in Calibri is
anachronistic.\cite{w-vistadate,w-segoe-default,w-cleartype,w-calibri}
\item \textbf{2009-10-22 --- Windows 7.} Segoe UI 9\,pt retained (Light/Semibold/Symbol
weights added).\cite{w-7date,w-segoe-default}
\item \textbf{2012-10-26 --- Windows 8.} Segoe UI ships a stylistic redraw (dated
October 2011), with a Semilight weight; Black and Emoji additions follow in 8.1
(2013-10-17).\cite{w-8date,w-segoe}
\item \textbf{2015-07-29 --- Windows 10.} Segoe UI retained. Structurally important
for provenance: \textbf{many script-specific fonts move into optional
Feature-on-Demand (FOD) packages}, installed automatically only when the associated
language is enabled---presence of a font on a Windows 10/11 system no longer implies
the OS installed it unconditionally.\cite{w-10date,w-fod}
\item \textbf{2021-10-05 --- Windows 11.} \textbf{Segoe UI Variable} (weight axis plus
an optical-size axis) becomes the default UI face; the Windows 11 font list marks
Segoe UI Variable, Segoe Fluent Icons, Cascadia Code, and Cascadia Mono ``Added in
Windows 11''.\cite{w-11list,w-11design}
\item \textbf{2023-07-13 --- Aptos announced as the Microsoft 365 default} (Steve
Matteson's ``Bierstadt'' renamed), replacing Calibri across Word, Excel, PowerPoint,
Outlook. Distribution boundary, precisely: Aptos reaches users (a) inside Microsoft
365 apps via the Office Font Service (``cloud fonts'') and (b) as a free manual
download from the Microsoft Download Center (v4.40 published 2024-06-04) for use
outside Office; it is \textbf{not} bundled as a Windows operating-system font and is
absent from Microsoft's Windows 11 font
list.\cite{w-aptos,w-aptos-2,w-aptos-dl,w-aptos-dl-2,w-11list}
\end{itemize}

Microsoft publishes authoritative per-version ``Fonts that ship with Windows'' lists
for Windows 7 through 11; no such pages exist for Vista or XP (retired URLs return
404), for which Wikipedia's typeface list with its ``Included from'' column is the
best available consolidated record.\cite{w-lists,w-typelist}

\subsection{Microsoft Windows --- CJK script lanes}\label{sec:win-cjk}

Figure~\ref{fig:win-cjk} traces the Windows CJK default UI fonts by script lane.

\begin{figure*}[t]
\centering
\begin{fptimeline}
\fpevent{1993}{Windows 3.1J---MS Gothic and MS Mincho by Ricoh licensed as the Japanese system fonts}
\fpevent{1998}{TC Windows 98---PMingLiU (New MingLiU) joins DynaLab MingLiU}
\fpevent{2000}{Windows 2000 (Feb 17)---one worldwide binary, full CJK set on all media (installable, not default-installed)---SimSun (SC), PMingLiU (TC), MS UI Gothic (JP), Gulim (KR) as locale UI fonts}
\fpevent{2007}{Windows Vista (Jan 30)---the four-lane ClearType switch---Microsoft YaHei (SC), Microsoft JhengHei (TC), Meiryo (JP), Malgun Gothic (KR); HKSCS-2004 native; ExtB fonts added}
\fpevent{2013}{Windows 8.1 (Oct 17)---Yu Gothic and Yu Mincho bundled (Jiyukobo); Meiryo UI still the JP UI font}
\fpevent{2015}{Windows 10 (Jul 29)---JP UI font becomes Yu Gothic UI; legacy CJK families move to Feature-on-Demand; DengXian ships as SC FOD font}
\fpevent{2021}{Windows 11 (Oct 5)---YaHei UI, JhengHei UI, Yu Gothic UI, Malgun Gothic all retained; Segoe UI Variable adds no CJK coverage}
\end{fptimeline}
\caption{Windows CJK default UI fonts by script lane.}
\label{fig:win-cjk}
\end{figure*}

\paragraph{Simplified Chinese (SC).} The ZhongYi (Beijing ZhongYi Electronics) faces
\textbf{SimSun 宋体} and \textbf{SimHei 黑体} anchor the lane from the regional Windows
3.x/95 editions (pre-2000 per-edition ship dates rest on secondary sources;
Microsoft's supply tables begin at Windows 2000, where SimSun/NSimSun v2.11 first
appear). SimSun was the \textbf{default SC UI font on Windows 2000/XP}. On
\textbf{Vista (2007-01-30)}, \textbf{Microsoft YaHei 微软雅黑}---glyphs licensed from
Beijing Founder Electronics, ClearType-optimized---became the default SC UI font;
SimSun remained bundled and \textbf{SimSun-ExtB} (CJK Extension B) was added.
\textbf{DengXian 等线} ships from Windows 10 as an optional Feature-on-Demand font (it
is also Office's Chinese default from Office 2016, \S\ref{sec:msoffice}). Windows 11
retains YaHei UI as the SC UI binding, since Segoe UI Variable covers only
Latin/Greek/Cyrillic.\cite{cw-simsun,cw-simsun-2,cw-yahei,cw-w10list,cw-w11,cw-w11-2}

\paragraph{Traditional Chinese (TC).} The DynaLab (later DynaComware, 華康) Ming
face---raster in TC Windows 3.0, TrueType from 3.1, renamed \textbf{MingLiU 細明體} at
v2.00 (``U'' for Unicode)---was the TC UI face from Windows 3.0 through XP, joined by
\textbf{DFKai-SB 標楷體} (kai) and, from TC Windows 98, the proportional \textbf{PMingLiU
新細明體} (default TC UI font on 2000/XP). On \textbf{Vista}, \textbf{Microsoft JhengHei
微軟正黑體} became the default TC UI font. Its designer is \textbf{China Type Design
Limited} of Hong Kong (a Monotype company)---not DynaComware, a frequent
misattribution. MingLiU-ExtB/PMingLiU-ExtB arrived with Vista; the whole legacy
Ming/Kai set moved to Feature-on-Demand in Windows 10. JhengHei UI remains the TC
binding in Windows 11.\cite{cw-mingliu,cw-mingliu-2,cw-mingliu-3,cw-kaplan,cw-jhenghei,cw-w10list}

\paragraph{Traditional Chinese --- Hong Kong (HKSCS).} Hong Kong's supplementary
character sets (GCCS 1995; HKSCS-1999/-2001/-2004/-2008/-2016) reached Windows first
as \textbf{patches} for Windows 98/NT4/2000/XP that enabled a hidden code page 951 and
patched MingLiU. \textbf{Vista} brought native, Unicode-based \textbf{HKSCS-2004}
support with dedicated \textbf{MingLiU\_HKSCS} and \textbf{MingLiU\_HKSCS-ExtB} fonts,
which remain the bundled HK glyph carriers (v7.03, in the TC Feature-on-Demand set)
through Windows 10/11.\cite{cw-hkscs,cw-hkscs-2}

\paragraph{Japanese (JP).} \textbf{MS Gothic} was ``designed by Ricoh Co. Ltd. and
licensed as the default system font for Windows 3.1'' (1993), with serif companion
\textbf{MS Mincho}; \textbf{MS UI Gothic} (from v2.30, Windows 2000) carried the JP UI
through XP. On \textbf{Vista}, \textbf{Meiryo \ja{メイリオ}}---designed by C\&G Inc.
(Satoru Akamoto, Takeharu Suzuki, Yukiko Ueda) with \textbf{Eiichi K\=ono} and
\textbf{Matthew Carter}, JIS X 0213:2004-compliant, part of the ClearType
collection---became the default JP system font; \textbf{Meiryo UI} followed in Windows
7. \textbf{Windows 8.1 (2013-10-17)} bundled Jiyukobo's \textbf{Yu Gothic
\ja{游ゴシック}} and \textbf{Yu Mincho \ja{游明朝}} while keeping Meiryo UI as the UI
face; the switch of the JP UI font to \textbf{Yu Gothic UI} came with \textbf{Windows
10 (2015-07-29)}, at which point Meiryo and MS Gothic/Mincho moved to the JP
supplemental FOD set. \textbf{UD Digi Kyokasho} (universal-design textbook faces)
arrived in Windows 10 version
1709.\cite{cw-msgothic,cw-meiryo,cw-meiryo-2,cw-yugothic,cw-fontsupport}

\paragraph{Korean (KR).} HanYang I\&C's \textbf{Gulim \ko{굴림}/GulimChe} (gothic) and
\textbf{Batang \ko{바탕}/BatangChe} (myeongjo) anchor the lane from 1990s Korean
Windows (Microsoft-documented from Windows 2000, v2.20); \textbf{Gulim} was the KR UI
font on 2000/XP. On \textbf{Vista}, \textbf{Malgun Gothic \ko{맑은 고딕}}---Sandoll
Communications (designers Kyoung-bae Lee and Daekwon Kim), Latin glyphs derived from
Segoe UI---became the default KR UI font and has remained so through Windows 11 (hanja
added in Windows 8 via China Type Design; Semilight weight later); Gulim/Batang/Dotum
moved to the KR FOD set in Windows 10.\cite{cw-gulim,cw-gulim-2,cw-malgun,cw-malgun-2}

\subsection{Apple macOS --- Latin lane}\label{sec:mac-latin}

Figure~\ref{fig:mac-latin} traces the classic Mac OS and macOS default system font.

\begin{figure*}[t]
\centering
\begin{fptimeline}
\fpevent{1984}{System 1 (Jan 24)---Chicago by Susan Kare, with the bitmap city fonts (Geneva, Monaco, New York and others)}
\fpevent{1991}{System 7 (May 13)---TrueType introduced}
\fpevent{1997}{Mac OS 8 (Jul 26)---Charcoal (Font Bureau, 1995) replaces Chicago as default; system font user-selectable}
\fpevent{2000}{Mac OS X Public Beta (Sep 13)---Lucida Grande debuts with Aqua}
\fpevent{2001}{Mac OS X 10.0 (Mar 24)---Lucida Grande confirmed as the OS X UI font}
\fpevent{2014}{OS X 10.10 Yosemite (Oct 16)---Helvetica Neue replaces Lucida Grande}
\fpevent{2015}{OS X 10.11 El Capitan (Sep 30)---San Francisco replaces Helvetica Neue; Apple's first in-house system typeface in two decades}
\end{fptimeline}
\caption{Classic Mac OS and macOS default system font (Latin lane).}
\label{fig:mac-latin}
\end{figure*}

\paragraph{Event register.}
\begin{itemize}[leftmargin=1.2em,itemsep=2pt]
\item \textbf{1984-01-24 --- System 1 (Macintosh 128K).} Default system font
\textbf{Chicago} (Susan Kare; 12\,pt bitmap), used for menus, dialogs and titles, and
retained as the system font through System 7.6. The original bundle is the bitmap
``city'' set---Chicago, Geneva, Monaco, New York, Venice, London, Athens, San
Francisco (the 1984 ransom-note face, unrelated to the 2015 system font), Toronto,
Cairo---largely Kare's work; her first-person account of the city-naming convention
survives on Folklore.org. The launch date is the 1984-01-24 shareholders'-meeting
unveiling, anchored institutionally by the Computer History
Museum.\cite{m-chicago,m-chicago-2,m-kare,m-kare-2,m-chm84,m-chm84-2}
\item \textbf{1991-05-13 --- System 7.} \textbf{TrueType} introduced---Apple's outline
font format, later licensed to Microsoft---freeing the Mac from fixed-size bitmaps (a
TrueType Chicago was cut to match the bitmap metrics). Date corroborated by
AppleInsider's dated retrospective; contemporaneous 1991 primary not
captured.\cite{m-truetype,m-sys7,m-sys7-2}
\item \textbf{1997-07-26 --- Mac OS 8.} Default system font changes to
\textbf{Charcoal} (designed 1995 at Font Bureau, based on Chicago), introduced with
the Platinum appearance; the Appearance control panel made the system font
user-selectable, with Chicago still available. The 1997-07-26 ship date rests on
secondary sources; Apple's 1997 press release survives only in Wayback snapshots that
could not be re-read during this research (archive.org rate limiting)---flagged
accordingly.\cite{m-charcoal,m-charcoal-2,m-os8,m-os8-2}
\item \textbf{2000-09-13 / 2001-03-24 --- Mac OS X Public Beta (``Kodiak''), then 10.0
(``Cheetah'').} \textbf{Lucida Grande} (Bigelow \& Holmes' Lucida family) becomes the
UI font of the Aqua interface, first in the Public Beta and then in 10.0, and remains
so until 2014. Apple's 2001-03-21 press release ``Mac OS X Hits Stores This Weekend''
fixes the 10.0 retail date (``beginning this Saturday, March 24'') and---a
bundled-font primary in its own right---advertises ``more than \$1{,}000 of the best
fonts'', naming Baskerville, Zapfino, Futura and Optima, plus ``the highest-quality
Japanese fonts'' (\S\ref{sec:apple-cjk}).\cite{m-lucida,m-pubbeta,m-osxpr}
\item \textbf{2014-10-16 --- OS X 10.10 Yosemite.} Default system typeface changes from
Lucida Grande to \textbf{Helvetica Neue}, aligning the Mac with iOS; the GA date is
fixed by Apple's press release. The font-switch attribution itself rests on the
encyclopedic record and John Siracusa's Ars Technica review (the review's typography
section could not be re-fetched during this research; Ars blocks automated
readers).\cite{m-yosemite-pr,m-yosemite-font,m-yosemite-font-2}
\item \textbf{2015-09-30 --- OS X 10.11 El Capitan.} Default system font changes from
Helvetica Neue to \textbf{San Francisco}, Apple's first in-house system typeface since
the 1980s bitmap era, introduced across platforms (iOS 9: 2015-09-16,
\S\ref{sec:apple-cjk}; watchOS earlier, 2015-04-24). Primary anchors: Apple's
availability press release and the WWDC 2015 session ``Introducing the New System
Fonts'' (Antonio Cavedoni).\cite{m-elcap-pr,m-wwdc804,m-sf}
\end{itemize}

\paragraph{San Francisco family evolution (bundled-font milestones).} SF Compact
shipped first on Apple Watch (watchOS, April 2015); \textbf{SF Mono} appeared at WWDC
2016 (Terminal/Xcode); the downloadable family was renamed from ``SF UI'' to
\textbf{SF Pro} in the WWDC 2017 timeframe---supported by the SF Pro font files' own
copyright string (``\copyright{} 2015--2017 Apple Inc.'') and the documented naming
history; the before/after Wayback snapshots of Apple's developer fonts page
(2016-06-18 vs 2017-06-12) exist per the availability API but their bodies could not
be re-read during this research, so the 2017 rebrand year is stated with that caveat.
\textbf{New York}, the serif companion (first seen as ``SF Serif'' in Apple Books, iOS
12), was released on the Apple Developer site on \textbf{2019-06-03}. Variable-font
builds of SF arrived at WWDC 2020, and the \textbf{SF script extensions} (SF Arabic at
WWDC 2021, later SF Armenian, SF Georgian, SF Hebrew) extend the family beyond
Latin.\cite{m-sf,m-sfpro,m-ny,m-devfonts}

\paragraph{Bundled-font inventories and removals.} Apple publishes a per-release
``Fonts included with macOS \textless name\textgreater'' support document (e.g.,
article 120414 for macOS Sequoia), distinguishing pre-installed fonts,
download-on-demand fonts (dimmed in Font Book), and document-support-only fonts---the
authoritative per-version inventory. Since the Sierra era many legacy families
(including the classic TC faces, \S\ref{sec:apple-cjk}, and Osaka) are
download-on-demand rather than pre-installed.\cite{m-sequoia-fonts}

\subsection{Apple --- CJK lanes (macOS and iOS)}\label{sec:apple-cjk}

Figure~\ref{fig:apple-cjk} traces the Apple CJK system fonts by script lane.

\begin{figure*}[t]
\centering
\begin{fptimeline}
\fpevent{1986}{KanjiTalk era---Osaka as the Japanese system face}
\fpevent{Mid-1990s}{Chinese Language Kit (version dates approximate)---Apple LiSung and Apple LiGothic (DynaLab) for TC; Song, Hei, Kai, Beijing for SC}
\fpevent{2000}{Apple announces Hiragino licensing (Feb 16, Macworld Tokyo)---six OpenType faces, 17,000 characters each}
\fpevent{2001}{Mac OS X 10.0 (Mar 24)---Hiragino ships; Kaku Gothic Pro is the JP UI face}
\fpevent{2002}{Mac OS X 10.2 Jaguar---STHeiti (SinoType) becomes the SC UI font}
\fpevent{2003}{Mac OS X 10.3 Panther---LiHei Pro and LiSong Pro (DynaComware, Big-5E with HKSCS) join; LiHei Pro is the TC UI face}
\fpevent{2009}{Mac OS X 10.6 Snow Leopard (Aug 28)---Hiragino Sans GB added for SC}
\fpevent{2011}{iOS 5---Apple SD Gothic Neo (Sandoll) replaces AppleGothic as the KR UI font on iOS}
\fpevent{2012}{OS X 10.8 Mountain Lion---Apple SD Gothic Neo becomes the KR UI font on the Mac}
\fpevent{2015}{iOS 9 (Sep 16) and OS X 10.11 (Sep 30)---PingFang SC, TC and HK become the Chinese UI fonts; Hiragino rebranded Hiragino Sans}
\end{fptimeline}
\caption{Apple CJK system fonts by script lane (classic Mac OS, macOS, iOS).}
\label{fig:apple-cjk}
\end{figure*}

\paragraph{Japanese.} The KanjiTalk-era system face was \textbf{Osaka} (bitmap;
TrueType by the System 7/KanjiTalk 7 era---classic-era dates approximate). The modern
era begins with Apple's \textbf{2000-02-16 Macworld Tokyo announcement}---still live on
Apple Newsroom---of an agreement with Dainippon Screen (fonts designed by
\textbf{Jiyukobo}) to bundle \textbf{Hiragino}: ``six typefaces, 17,000 characters
each, in OpenType'', the JIS X 0213 expanded set. Hiragino shipped with Mac OS X 10.0
(2001-03-24), and the Hiragino gothic faces have been the Japanese UI lineage on Mac
OS X since---initially as \textbf{Hiragino Kaku Gothic Pro}, represented since the El
Capitan era by the reorganized \textbf{Hiragino Sans} family (weights W0--W9). Osaka
persists as a download-on-demand legacy
font.\cite{ca-hiragino-pr,ca-hiragino,ca-sierra}

\paragraph{Simplified Chinese.} Mac OS X bundled Changzhou SinoType's ST faces;
\textbf{STHeiti became the SC UI font in 10.2 Jaguar} (2002)---not 10.0, a common
conflation---with the full outline family (STXihei, STSong, STKai, STFangsong) present
by 10.3. \textbf{Hiragino Sans GB} (an SC adaptation with Beijing Hanyi) was added in
10.6 Snow Leopard (2009-08-28). The ST families were regrouped as Heiti SC / Songti SC
/ Kaiti SC collections in the 10.7/10.8 era.\cite{ca-chinesemac,ca-cjklist}

\paragraph{Traditional Chinese and Hong Kong.} The Chinese Language Kit era
($\sim$1995) brought DynaLab's \textbf{Apple LiSung 儷宋} and \textbf{Apple LiGothic
儷中黑} plus BiauKai 標楷. 10.3 Panther added \textbf{LiHei Pro 儷黑 Pro} and
\textbf{LiSong Pro 儷宋 Pro}---Big-5E character set \textbf{including HKSCS}, making them
the HK carriers---with LiHei Pro as the TC UI face through 2015. The legacy TC faces
were demoted to download-on-demand around the Sierra era (not removed in Catalina, a
common misreport).\cite{ca-chinesemac,ca-sierra}

\paragraph{The 2015 PingFang event.} With iOS 9 (2015-09-16) and OS X 10.11 El Capitan
(2015-09-30), \textbf{PingFang 蘋方} became the Chinese UI family---produced with
DynaComware, six weights, in three regional variants: \textbf{PingFang SC, PingFang
TC, and PingFang HK}. This is Apple's first dedicated Hong Kong UI lane (previously HK
support rode inside Taiwan-TC fonts via HKSCS), and the event replaced STHeiti/Heiti
SC/TC and LiHei Pro simultaneously.\cite{ca-pingfang,ca-pingfang-2,ca-chinesemac11}

\paragraph{Korean.} \textbf{AppleGothic} carried the KR UI from the Korean Language Kit
era through 2011--2012, when Sandoll's \textbf{Apple SD Gothic Neo} replaced it---first
on \textbf{iOS 5 (2011)}, then on \textbf{OS X 10.8 Mountain Lion (2012)} with nine
weights; it remains bundled (macOS Sequoia ships v19.x). The iOS 5.0-vs-5.1 point and
the Mountain Lion attribution rest on foundry/community records rather than an Apple
statement---flagged accordingly.\cite{ca-korean,ca-korean-2}

\paragraph{iOS bundled-font history (Latin note in passing).} The original iPhone
(2007) shipped a curated subset of Mac faces---Gruber's contemporaneous PDF specimen,
still live, classifies 9 full families installed (Arial, Courier New, Georgia,
Helvetica, Marker Felt, Times New Roman, Trebuchet MS, Verdana, Zapfino), 2 partial, 3
name-substituted (\textbf{Helvetica Neue substituted by Helvetica}---Neue was not
actually installed until the iPhone 4 era), and 27 Mac families absent
(\S\ref{sec:ios-latin}).\cite{ca-gruber,ca-gruber-2}

\subsection{Apple iOS --- Latin lane}\label{sec:ios-latin}

Figure~\ref{fig:ios-latin} traces the iOS default system font.

\begin{figure}[t]
\centering
\begin{fptimeline}[0.72\columnwidth]
\fpevent{2007}{iPhone OS 1.0 (Jun 29)---Helvetica as system font; Helvetica Neue not installed, only name-substituted}
\fpevent{2010}{iPhone 4 with iOS 4 (Jun 21 software, Jun 24 device)---Helvetica Neue on Retina hardware; non-Retina devices keep Helvetica}
\fpevent{2013}{iOS 7 (Sep 18)---flat redesign leans on Helvetica Neue Light and Ultralight; beta 3 walks small text back to Regular}
\fpevent{2015}{iOS 9 (Sep 16)---San Francisco replaces Helvetica Neue}
\end{fptimeline}
\caption{iOS default system font.}
\label{fig:ios-latin}
\end{figure}

\paragraph{Event register.}
\begin{itemize}[leftmargin=1.2em,itemsep=2pt]
\item \textbf{2007-06-29 --- iPhone OS 1.0.} System font \textbf{Helvetica}; Notes uses
Marker Felt. Helvetica Neue was \emph{not installed}---pages requesting it silently got
Helvetica (Gruber's specimen, \S\ref{sec:apple-cjk}). The full inventory: 9 full
families, 2 partial, 3 substituted names, 27 absent.\cite{ca-gruber,ca-gruber-2,i-gruber07}
\item \textbf{2010-06-21 / 2010-06-24 --- iOS 4 and iPhone 4.} \textbf{Helvetica Neue
becomes the system font on Retina hardware.} This change is device-driven, not an
OS-version change: the same iOS 4 on non-Retina devices (3GS and earlier) kept
Helvetica---Gruber's contemporaneous review is explicit, attributing the split to
hinting quality on low-density screens. Apple's iPhone 4 press release carries both
dates (iOS 4 free update June 21; iPhone 4 on sale June 24) and the Retina
display.\cite{i-gruber4,i-iphone4pr}
\item \textbf{2013-09-18 --- iOS 7.} The flat redesign makes \textbf{Helvetica Neue}
carry the interface system-wide, leaning on \textbf{Light/Ultralight} weights, with
Dynamic Type introduced. In the beta cycle, betas 1--2 set most UI text in Helvetica
Neue Light and were widely criticized; \textbf{beta 3 (2013-07-08) walked small text
back to Regular}---a documented within-cycle default correction. GA date per Apple's
press release.\cite{i-ios7pr,i-typographica,i-fortune}
\item \textbf{2015-09-16 --- iOS 9.} \textbf{San Francisco} replaces Helvetica Neue as
the system font, in lockstep with El Capitan (\S\ref{sec:mac-latin}); SF had debuted on
Apple Watch on 2015-04-24 and been available to developers since 2014-11-18. GA date
per Apple's press release.\cite{i-ios9pr,m-sf}
\item \textbf{2019-09-19 --- iOS 13 (provenance-relevant API change).} Third-party
\textbf{custom font installation} becomes supported system-wide---fonts distributed
through App Store apps, registered via \code{CTFontManagerRegisterFontURLs}, managed in
Settings (WWDC 2019 session 227). After this point, presence of a non-Apple font on an
iOS device no longer implies jailbreak or configuration-profile
sideloading.\cite{i-ios13,i-ios13-2,i-ios13-3}
\end{itemize}

\subsection{Android}\label{sec:android}

Figure~\ref{fig:android} traces the Android default fonts (Latin and CJK).

\begin{figure*}[t]
\centering
\begin{fptimeline}
\fpevent{2007}{Droid family announced (Nov 12)---Ascender Corporation, designer Steve Matteson, for the Open Handset Alliance}
\fpevent{2008}{Android 1.0 (Sep 23)---Droid Sans default UI, Droid Serif and Droid Sans Mono bundled, Droid Sans Fallback for CJK}
\fpevent{2011}{Android 4.0 Ice Cream Sandwich (Oct 19 SDK)---Roboto replaces Droid Sans; MotoyaLMaru and MotoyaLCedar bundled for Japanese}
\fpevent{2014}{Roboto v2 published (Jul 16) with Material Design; Noto Sans CJK released (Jul 15); both ship in Android 5.0 Lollipop (Nov 12)}
\fpevent{2015}{Android 6.0 Marshmallow---legacy Droid Sans Fallback and Motoya files removed, Noto covers all CJK}
\fpevent{2017}{Android 8.0 Oreo (Aug 21)---Downloadable Fonts API and fonts-in-XML; Google Sans appears on Pixel skin (8.1)}
\fpevent{2022}{Roboto Flex variable font released to Google Fonts (May 5)---not adopted as the AOSP system default}
\end{fptimeline}
\caption{Android default fonts (Latin and CJK).}
\label{fig:android}
\end{figure*}

\paragraph{Event register (Latin lane).}
\begin{itemize}[leftmargin=1.2em,itemsep=2pt]
\item \textbf{2007-11-12 --- Droid family announced.} Ascender Corporation announces
the Droid collection (Droid Sans, Droid Serif, Droid Sans Mono), designed by
\textbf{Steve Matteson}, commissioned for the Open Handset Alliance's Android platform.
The original PRWeb release is dead; a Wayback snapshot exists (existence confirmed via
the availability API; body not re-read during this research due to archive.org rate
limiting), with content corroborated by the encyclopedic record.\cite{a-droid-pr,a-droid}
\item \textbf{2008-09-23 --- Android 1.0.} \textbf{Droid Sans} is the default UI font;
\textbf{Droid Sans Fallback}---Apache-licensed, $\sim$15,500 Han glyphs built by
component reference, following mainland-PRC preferred glyph shapes---is the single CJK
fallback (the root of later ``wrong-region glyph'' complaints in TC/JP contexts).
Primary date anchor: the Android Developers Blog post announcing the Android 1.0 SDK,
release 1.\cite{a-10sdk,a-droid,a-fallback}
\item \textbf{2011-10-19 --- Android 4.0 Ice Cream Sandwich (SDK; first device Galaxy
Nexus, Nov 2011).} \textbf{Roboto} (Christian Robertson) replaces Droid Sans as the
system font with the Holo design language. Note on sourcing: the Android 4.0 platform
announcement post itself does not mention Roboto---the release/date primary is the
platform post, while the Roboto attribution rests on the ICS-era Android Design
typography page (Wayback snapshot; body not re-read this research) and contemporaneous
press. Japanese: \textbf{MotoyaLMaru/MotoyaLCedar} (Motoya Co.) bundled, MotoyaLMaru
serving as the JP fallback ahead of Droid Sans Fallback (these later became Google
Fonts' Kosugi/Kosugi Maru).\cite{a-ics-sdk,a-ics-typo,a-wired,a-motoya,a-motoya-2}
\item \textbf{2014-07-16 --- Roboto v2 published.} Christian Robertson's redesigned
Roboto (``The New Roboto'', Google Developers Blog) accompanies Material Design
(announced at Google I/O 2014-06-25): rounder forms, new weights. It ships as the
system default with \textbf{Android 5.0 Lollipop (2014-11-12)}---publication and OS
shipment are distinct dated events.\cite{a-roboto2,a-roboto2-2,a-lollipop}
\item \textbf{2017-08-21 --- Android 8.0 Oreo.} \textbf{Downloadable Fonts API} and
fonts-as-XML-resources (API 26)---apps no longer need to bundle font files;
provenance-relevant because app typography stops implying APK contents.\cite{a-oreo}
\item \textbf{Pixel vs AOSP (standing distinction).} Google's proprietary \textbf{Product
Sans/Google Sans} appears on Pixel devices from Android 8.1 (Dec 2017) in branding
surfaces, with \textbf{Google Sans Text} spreading through Pixel system apps around
Android 12 (2021)---but these are Pixel skin overlays; \textbf{the AOSP default remains
Roboto throughout}, and \textbf{Roboto Flex} (the 12-axis variable font released to
Google Fonts on 2022-05-05) has \emph{not} been adopted as the AOSP system default.
Pixel-lane version attributions rest on secondary observation and are marked
UNCERTAIN.\cite{a-gsans,a-flex,a-flex-2}
\end{itemize}

\paragraph{CJK lane.}
\begin{itemize}[leftmargin=1.2em,itemsep=2pt]
\item \textbf{2014-07-15 --- Source Han Sans / Noto Sans CJK released.} The first
open-source Pan-CJK family---a joint Adobe (Source Han Sans) and Google (Noto Sans CJK)
release, identical fonts under two brands, with partner foundries SinoType (SC), Iwata
(JP), and Sandoll (KR); seven weights with distinct SC/TC/JP/KR regional
variants.\cite{a-shs,a-shs-2}
\item \textbf{2014-11-12 --- Android 5.0 Lollipop.} \textbf{Noto Sans CJK becomes
Android's CJK default}, replacing Droid Sans Fallback and delivering per-region glyph
correctness for the first time (\code{fonts.xml} with \code{zh-Hans}/\code{zh-Hant}/
\code{ja}/\code{ko} language tags; only the Regular weight integrated in
practice).\cite{a-shs,a-shs-2,a-shs75}
\item \textbf{2015 --- Android 6.0 Marshmallow.} Legacy \textbf{Droid Sans Fallback and
Motoya files removed} once Noto coverage is complete---corroborated by a
system-partition diff of Nexus builds 5.1.1 vs 6.0.0; from Android 7.0 region selection
moved to \code{ttcIndex} entries within a single Noto CJK collection
file.\cite{a-cjk-remove}
\item \textbf{2017-04-03 --- Source Han Serif / Noto Serif CJK released} (the serif
companion). Whether/when stock Android bundled a serif CJK by default is \textbf{not
established}---no AOSP primary found; treat serif-CJK presence on a device as app- or
vendor-supplied absent other evidence. Variable Noto CJK on stock Android is
community-attested for \textbf{Android 15+} only.\cite{a-shserif,a-cjk-vf,a-cjk-vf-2}
\end{itemize}

\subsection{Desktop Linux}\label{sec:linux}

Figure~\ref{fig:linux} traces desktop Linux default UI fonts.

\begin{figure*}[t]
\centering
\begin{fptimeline}
\fpevent{1999}{Arphic donates four Chinese TrueType fonts under the Arphic Public License}
\fpevent{2003}{Bitstream Vera released (Jan 22, GNOME Foundation agreement); Kochi incident deprecates the free Japanese defaults (Jun 27)}
\fpevent{2004}{DejaVu forks Vera---becomes the de facto distro default for years}
\fpevent{2007}{Liberation fonts (May 9, Red Hat and Ascender)---metric-compatible with Arial, Times New Roman, Courier New; WenQuanYi Zen Hei (Sep 15) for Chinese}
\fpevent{2010}{Ubuntu 10.10 (Oct 10)---Ubuntu Font Family (Dalton Maag) becomes the Ubuntu default}
\fpevent{2011}{GNOME 3.0 (Apr 6)---Cantarell (Dave Crossland) becomes the GNOME default}
\fpevent{2014}{KDE Plasma 5.0 (Jul 15)---Breeze theme with Noto Sans as default}
\fpevent{2016}{Noto unified release (Oct 6); Ubuntu 16.04 switches Chinese to Noto Sans CJK}
\fpevent{2018}{Fedora 29 switches all CJK defaults to Noto; Ubuntu 18.04 switches Japanese to Noto}
\fpevent{2023}{Debian bookworm era---fontconfig 2.14 switches generic Latin defaults from DejaVu to Noto}
\fpevent{2025}{GNOME 48 (Mar 19)---Adwaita Sans (an Inter fork) replaces Cantarell; Adwaita Mono (Iosevka) replaces Source Code Pro}
\end{fptimeline}
\caption{Desktop Linux default UI fonts---foundational packages and desktop switches.}
\label{fig:linux}
\end{figure*}

\paragraph{Foundational packages.}
\begin{itemize}[leftmargin=1.2em,itemsep=2pt]
\item \textbf{1999 --- Arphic Public License release.} Taiwan's Arphic Technology (文鼎)
donates four TrueType fonts---AR PL Mingti2L Big5, AR PL KaitiM Big5 (TC), AR PL
SungtiL GB, AR PL KaitiM GB (SC)---under the copyleft \textbf{Arphic Public License}
(not the GPL, a common miscitation). These are the foundational free Chinese fonts;
firefly's bitmap-patched derivatives became the \textbf{AR PL UMing/UKai} collections
that long served as Debian/Ubuntu Chinese defaults.\cite{l-arphic,l-arphic-2,l-uming}
\item \textbf{2003-01-22 --- Bitstream Vera.} The GNOME Foundation announces the
agreement with Bitstream to release the ten-font Vera family (Jim Lyles) under a
permissive license---the de facto Linux desktop faces of the mid-2000s.\cite{l-vera}
\item \textbf{2004 --- DejaVu.} \v{S}t\v{e}p\'an Roh forks Vera 1.10 to extend coverage;
DejaVu becomes the default sans/serif/mono across Fedora, Debian, Ubuntu and openSUSE
for roughly 2005--2011 and remains a fontconfig fallback staple (final release 2.37,
2016).\cite{l-dejavu}
\item \textbf{2007-05-09 --- Liberation fonts.} Red Hat (with Ascender; design Steve
Matteson) releases Liberation Sans/Serif/Mono, \textbf{metrically identical to Arial,
Times New Roman, and Courier New}---``when substituted for the Microsoft fonts, a line
of text is identically displayed'' per Red Hat's announcement. Relicensed SIL OFL
(v2.0, 2012, rebased on Croscore). Default document faces in Fedora and LibreOffice
(\S\ref{sec:staroffice}).\cite{l-liberation,l-liberation-2}
\item \textbf{2016-10-06 --- Noto unified release.} Google/Monotype's ``no more tofu''
family publicly launches as a unified whole (components shipped earlier---Android used
Noto CJK from 2014, \S\ref{sec:android}).\cite{l-noto}
\end{itemize}

\paragraph{Desktop/distro default switches.}
\begin{itemize}[leftmargin=1.2em,itemsep=2pt]
\item \textbf{2010-10-10 --- Ubuntu 10.10.} The \textbf{Ubuntu Font Family} (Dalton
Maag, funded by Canonical, Ubuntu Font Licence 1.0) becomes the default UI font,
replacing DejaVu Sans; Launchpad confirms the 10.10 default installation and
authorship.\cite{l-ubuntufont}
\item \textbf{2011-04-06 --- GNOME 3.0.} \textbf{Cantarell} (Dave Crossland) becomes
GNOME's default UI font---GNOME's own announcement: ``the GNOME desktop will use
Cantarell by default.'' Fedora inherits it from Fedora 15 (May 2011) as the Workstation
default.\cite{l-cantarell,l-cantarell-2}
\item \textbf{2014-07-15 --- KDE Plasma 5.0.} The Breeze theme makes \textbf{Noto Sans}
the Plasma default. Precision note: KDE's own 5.0 announcement (primary for the date
and the Breeze theme) does not name the font; the Noto attribution rests on distro
configuration records and community documentation, and is labeled
accordingly.\cite{l-plasma5}
\item \textbf{2022--2023 --- Debian's generic aliases move to Noto.} fontconfig 2.14
changed the generic Latin defaults (sans/serif) from DejaVu to Noto in
\code{60-latin.conf}, landing in the Debian bookworm era; after objections Debian
reverted the \emph{monospace} alias to DejaVu Sans Mono (fontconfig 2.14.2-4). This is
the fontconfig alias layer, distinct from any desktop's UI font.\cite{l-debian}
\item \textbf{2025-03-19 --- GNOME 48.} \textbf{Adwaita Sans}---a customized build of
Rasmus Andersson's \textbf{Inter}---replaces Cantarell as GNOME's UI font, and
\textbf{Adwaita Mono} (an Iosevka build) replaces Source Code Pro: the first GNOME
UI-font change since 3.0. Fedora 42 ships it via GNOME 48. GNOME's release notes
describe the new fonts; the replaced-predecessor framing is corroborated by Phoronix
and the GNOME development blog.\cite{l-gnome48,l-gnome48-2,l-gnome48-3}
\item \textbf{RHEL note.} Red Hat Display/Text (2019, MCKL) are \textbf{brand fonts},
not the RHEL desktop UI default---RHEL's GNOME desktop follows the GNOME lineage above.
Asserting ``RHEL defaults to Red Hat Display'' would be incorrect.\cite{l-redhat,l-redhat-2}
\end{itemize}

\paragraph{CJK lanes on Linux.}
\begin{itemize}[leftmargin=1.2em,itemsep=2pt]
\item \textbf{Japanese.} The free-font lineage runs Kochi (deprecated after the
2003-06-27 glyph-copying disclosure against TypeBank's Mincho-M) $\rightarrow$ Sazanami
(2004, interim) $\rightarrow$ \textbf{VL Gothic} (Fedora's JP default from Fedora 9,
2008) and \textbf{IPA fonts} (open-sourced under the OSI-approved IPA Font License,
2009-04-01) $\rightarrow$ \textbf{Takao} forks (Ubuntu JP default from 10.04, first
release 2010-02-15) $\rightarrow$ \textbf{Noto Sans CJK JP} (Ubuntu 18.04; Fedora
29).\cite{l-kochi,l-vlgothic,l-ipa,l-takao,l-noto-switch,l-noto-switch-2,l-noto-switch-3}
\item \textbf{Chinese.} Arphic/UMing/UKai (above) $\rightarrow$ \textbf{WenQuanYi Zen
Hei} (released 2007-09-15; project founded 2004 by Qianqian Fang), the first
open-source Chinese heiti, default for Simplified Chinese in the Fedora/Ubuntu
2007--2008 cycle (exact first-default releases UNCERTAIN) $\rightarrow$ \textbf{Noto
Sans CJK} (Ubuntu 16.04 for Chinese---with a famous fontconfig ordering bug that let
the JP variant win for zh users until reprioritized; Fedora 29
system-wide).\cite{l-wqy,l-noto-switch,l-noto-switch-2,l-noto-switch-3}
\item \textbf{Korean.} Baekmuk (1990s) $\rightarrow$ UnFonts (Ubuntu default $\sim$2007)
$\rightarrow$ \textbf{Nanum} (Naver, released 2008; Ubuntu's Korean default from 12.04)
$\rightarrow$ Noto CJK KR (Fedora 29).\cite{l-nanum,l-nanum-2,l-noto-switch,l-noto-switch-2,l-noto-switch-3}
\item \textbf{The Fedora 29 consolidation (2018).} Fedora's Change
``CJKDefaultFontsToNoto'' (owners Akira Tagoh, Peng Wu) formally switched \textbf{all}
CJK defaults---VL Gothic (JP), NanumGothic (KR), and the Source Han variants
(SC/TC)---to Noto CJK. Negative finding worth recording: the widely-repeated claim that
Fedora 21 (2014) made Noto the CJK default has \textbf{no primary Fedora Change behind
it}; the documented system-wide switch is
F29.\cite{l-noto-switch,l-noto-switch-2,l-noto-switch-3}
\end{itemize}

% === §4 Office suites ========================================================
\section{Office Suites}\label{sec:office}

\subsection{Microsoft Office}\label{sec:msoffice}

Figure~\ref{fig:msoffice} traces the Microsoft Office default document fonts across
Latin and CJK edition lanes.

\begin{figure*}[t]
\centering
\begin{fptimeline}
\fpevent{1989}{Word for Windows era---Times New Roman 10\,pt (Excel Arial 10)}
\fpevent{1999}{Office 2000 (Jun 7)---Word default becomes Times New Roman 12\,pt}
\fpevent{2007}{Office 2007 (Jan 30 retail)---Calibri 11\,pt replaces Times New Roman in Word and Arial in Excel, PowerPoint, Outlook}
\fpevent{2013}{Office 2013---Calibri body with Calibri Light headings}
\fpevent{2015}{Office 2016 (Sep 22)---Japanese Word body default becomes Yu Mincho (Excel and PowerPoint follow Yu Gothic theme fonts); Simplified Chinese edition switches to DengXian; Traditional Chinese stays PMingLiU}
\fpevent{2021}{Beyond Calibri (Apr 28)---five candidate successors announced}
\fpevent{2023}{Aptos announced as new default (Jul 13); M365 rollout through early 2024}
\fpevent{2024}{Office LTSC 2024 and Office 2024 (Sep 16 and Oct 1)---first perpetual releases defaulting to Aptos}
\end{fptimeline}
\caption{Microsoft Office default document fonts (Latin and CJK edition lanes).}
\label{fig:msoffice}
\end{figure*}

\paragraph{Latin lane.}
\begin{itemize}[leftmargin=1.2em,itemsep=2pt]
\item \textbf{Word for Windows through Word 97:} Normal-template default \textbf{Times
New Roman 10\,pt}; Excel \textbf{Arial 10}. \textbf{Office 2000 (GA 1999-06-07)} raised
Word's default to \textbf{Times New Roman 12\,pt}---at Word 2000, not 2002/2003 as
commonly misremembered---persisting through Word 2003.\cite{o-barnhill,o-barnhill-2}
\item \textbf{Office 2007 (volume-license 2006-11-30; retail GA 2007-01-30):}
\textbf{Calibri 11\,pt} replaces Times New Roman as the Word default and Arial as the
Excel/PowerPoint/Outlook default, with Cambria as the serif heading companion---the
arrival of the ClearType Font Collection (\S\ref{sec:win-latin}). The change was
championed by Word program management (Joe Friend) for screen readability. The
``Fontgate'' affair (2017) turns on the \emph{earlier} of Calibri's two public
boundaries: the disputed documents were represented as executed in February 2006,
before even the Office 2007 Beta 2 public distribution of 2006-05-23---let alone the
2007-01-30 general availability.\cite{w-calibri,o-friend,o-fontgate}
\item \textbf{Office 2010--2021:} Calibri retained; Office 2013 introduced the
Calibri-plus-Calibri-Light theme (``Office Theme 2013--2022'').\cite{o-theme}
\item \textbf{2021-04-28 --- ``Beyond Calibri'':} Microsoft announces five commissioned
candidates (Bierstadt, Tenorite, Skeena, Seaford, Grandview), all added to the
Microsoft 365 font picker. \textbf{2023-07-13:} Bierstadt, renamed \textbf{Aptos}
(Steve Matteson), is announced as the new default; staged rollout makes it the
Microsoft 365 default through late 2023--early 2024, and \textbf{Office LTSC 2024 /
Office 2024} (2024-09-16 / 2024-10-01) are the first perpetual releases defaulting to
Aptos.\cite{o-beyond,w-aptos,w-aptos-2,o-2024}
\item \textbf{Cloud fonts:} Microsoft 365's on-demand ``cloud fonts'' arrived in the
2016--2017 window (exact launch date UNCERTAIN against Microsoft documentation)---
provenance-relevant because subscriber documents may use fonts never locally
installed.\cite{o-cloud}
\item \textbf{Office for Mac:} Office 2008 for Mac reportedly defaulted to Cambria (a
Mac-vs-Windows inconsistency), aligning to Calibri by Office 2016 for Mac---[secondary],
with per-app details UNCERTAIN.\cite{o-mac}
\end{itemize}

\paragraph{CJK edition lanes.} Office binds per-script default fonts in its theme
(\code{theme1.xml}, ISO 15924 script tags), so each language edition has its own
default lineage:
\begin{itemize}[leftmargin=1.2em,itemsep=2pt]
\item \textbf{Japanese:} \textbf{MS Mincho} (body) for decades $\rightarrow$
\textbf{Office 2016 for Windows (2015-09-22): Yu Mincho} body with \textbf{Yu Gothic}
headings (Excel: MS P Gothic $\rightarrow$ Yu Gothic). The Yu fonts ship in Windows
from 8.1 (\S\ref{sec:win-cjk}). JustSystems' Ichitaro made the equivalent switch in
2021 (\S\ref{sec:cjkoffice}).\cite{o-jp}
\item \textbf{Simplified Chinese:} \textbf{SimSun 宋体} $\rightarrow$ \textbf{DengXian
等线} as the Office 2016 SC-edition default. Precision matters here: the DengXian
\emph{font} first shipped in \textbf{Windows 10 version 1507 (2015-07-29)} as part of
the SC supplemental fonts; \textbf{Office 2016} adopted it as the SC default theme
font.\cite{o-dengxian}
\item \textbf{Traditional Chinese:} \textbf{PMingLiU 新細明體 --- no change.} The TC
edition kept PMingLiU as its default through current Microsoft 365; Microsoft JhengHei
(Vista-era ClearType sans, \S\ref{sec:win-cjk}) is available but is not the document
default. A TC document defaulting to PMingLiU alongside an SC document defaulting to
DengXian is thus itself a script-edition discriminator.\cite{o-tc}
\item \textbf{Korean:} Windows' Korean \emph{UI} default moved Gulim $\rightarrow$
Malgun Gothic at Vista (\S\ref{sec:win-cjk}), but \textbf{which Office release (if any)
switched the Korean document body default from Batang-era faces to Malgun Gothic is
UNCERTAIN}---no clean primary source found; not asserted.\cite{o-kr}
\item \textbf{Aptos-era theme bindings (verified by artifact inspection):} Aptos covers
Latin/Greek/Cyrillic (plus Vietnamese) only, so the Aptos theme resolves CJK text
through per-script \code{<a:font>} bindings. Direct inspection of the bundled
\code{Office Theme.thmx} $\rightarrow$ \code{theme1.xml} in an installed Microsoft Word
for Mac (Microsoft 365, version 16.111.1, observed 2026-07-24) shows: \textbf{majorFont}
(headings) latin = \emph{Aptos Display}, with Jpan = \ja{游ゴシック} Light (Yu Gothic
Light), Hang = \ko{맑은 고딕} (Malgun Gothic), Hans = 等线 Light (DengXian Light), Hant =
新細明體 (PMingLiU); \textbf{minorFont} (body) latin = \emph{Aptos}, with Jpan =
\ja{游ゴシック}, Hang = \ko{맑은 고딕}, Hans = 等线, Hant = 新細明體. Note the Japanese
nuance this exposes: the \emph{theme} body binding is \textbf{Yu Gothic}, while Word's
Japanese-edition document body default is documented as \textbf{Yu Mincho}
[secondary]---which implies a template-level override of the theme binding in ja-JP Word
(the ja-JP Normal template itself was not inspected here), whereas Excel/PowerPoint
follow the theme's Yu Gothic.\cite{o-aptos-cjk}
\end{itemize}

\subsection{StarOffice, OpenOffice.org, LibreOffice}\label{sec:staroffice}

Figure~\ref{fig:staroffice} traces the StarOffice, OpenOffice.org and LibreOffice
defaults.

\begin{figure*}[t]
\centering
\begin{fptimeline}
\fpevent{1999}{StarOffice 5.x (Sun)---Agfa Monotype metric clones bundled---Albany (Arial), Thorndale (Times New Roman), Cumberland (Courier New)}
\fpevent{2002}{OpenOffice.org 1.0 (May 1)---relies on system fonts; Writer default carried as Times New Roman by name}
\fpevent{2008}{OOo 2.4 (Mar)---Liberation fonts bundled}
\fpevent{2011}{LibreOffice 3.3 (Jan 25)---Writer default is Liberation Serif 12\,pt, Impress and Calc use Liberation Sans}
\fpevent{2013}{LibreOffice 4.1 (Jul 25)---Carlito and Caladea bundled as Calibri and Cambria metric substitutes}
\end{fptimeline}
\caption{StarOffice, OpenOffice.org and LibreOffice defaults.}
\label{fig:staroffice}
\end{figure*}

\begin{itemize}[leftmargin=1.2em,itemsep=2pt]
\item \textbf{StarOffice 5.x (Sun era, 1999--2000).} Bundled Agfa Monotype's
metric-compatible clones---\textbf{Albany} ($\triangleq$ Arial), \textbf{Thorndale}
($\triangleq$ Times New Roman), \textbf{Cumberland} ($\triangleq$ Courier New)---so
documents exchanged with Microsoft Office kept identical line breaks. The ``AMT''
suffix on enterprise systems is ``Agfa Monotype''. Precise first-bundling point release
UNCERTAIN; StarOffice 5.2 released 2000-06-20.\cite{s-staroffice,s-staroffice-2}
\item \textbf{OpenOffice.org.} 1.0 released 2002-05-01 (announcement 2002-04-30); the
proprietary AMT clones could not be redistributed, so OOo leaned on system fonts, with
the Writer default carried as ``Times New Roman'' \emph{by name} (substituted on
Linux)---the exact default string is UNCERTAIN against primary documentation. Bundling
history: Bitstream Vera (through 2.3) $\rightarrow$ \textbf{Liberation from 2.4}
(2008-03) $\rightarrow$ Gentium added at 3.2 (2010-02).\cite{s-ooo,s-ooo-2}
\item \textbf{LibreOffice.} Forked 2010-09-28; 3.3 released 2011-01-25. \textbf{Writer's
default has been Liberation Serif 12\,pt from 3.3 to the present} (Impress/Calc:
Liberation Sans)---no shipped change found across 3.x--25.x, itself a provenance
finding. \textbf{LibreOffice 4.1 (2013-07-25) began bundling Carlito and Caladea}
(Google's ``crosextra'' faces, metric-compatible with Calibri and Cambria), the
substitution pair that makes post-2007 \code{.docx} files lay out correctly without
Microsoft fonts---a Carlito rendering of Calibri-styled text indicates metric
substitution on a host lacking Calibri (LibreOffice bundles Carlito; so do many Linux
distributions), not uniquely a LibreOffice host. CJK: LibreOffice's per-locale default
lists (\code{VCL.xcu}) prefer Source Han Sans/Noto Sans CJK first, falling back to
platform faces (Microsoft YaHei/NSimSun on Windows, PingFang on macOS), so the effective
CJK default is host-dependent.\cite{s-lo,s-lo-2,s-lo41,s-lo41-2,s-lo41-3,s-lo-cjk}
\end{itemize}

\subsection{WPS Office, Google Docs, Apple iWork}\label{sec:wps}

\begin{itemize}[leftmargin=1.2em,itemsep=2pt]
\item \textbf{WPS Office (Kingsoft).} For Chinese documents WPS follows the SimSun
convention---behavioral testing shows 宋体 (SimSun) as the effective Chinese default and
missing-font fallback (STSong on macOS), matching Chinese-locale Word; its Latin default
is UNCERTAIN against vendor documentation. WPS has integrated third-party free fonts
(e.g., Xiaomi's MiSans) rather than shipping a proprietary clone set. It is widely
deployed in China, including a dedicated government ``Official Document Edition'' (2022);
market-share figures circulating for WPS (e.g., a $\sim$95\%-by-1993 claim) survive only
in tertiary compilations and are not asserted here.\cite{s-wps,s-wps-2}
\item \textbf{Google Docs.} Default body font \textbf{Arial 11} with 1.15 line spacing,
as documented for recent Google Docs; no record of a default-face change was found in
the sources searched, but the earliest-years (2006-era) default is unverified---the
continuous-since-2006 reading is UNCERTAIN. Google Fonts (launched May 2010) later fed
Docs' ``more fonts'' picker, expanding available---not default---fonts.\cite{s-gdocs,s-gdocs-2}
\item \textbf{Apple iWork Pages.} Recent Pages versions default the Blank template body
to \textbf{Helvetica Neue 11\,pt} per secondary documentation (Pages has no global
default font---templates' paragraph styles carry it); this was not version-verified
against a named Pages build during this research. Pre-2013 template defaults (Pages
'05--'09) are UNCERTAIN against Apple primary documentation and are not asserted
here.\cite{s-pages,s-pages-2}
\item \textbf{Historical one-liners, explicitly unverified.} WordPerfect for Windows
(commonly reported Times New Roman 12, earlier versions Courier) and Lotus Word Pro
(commonly reported Times New Roman)---both UNCERTAIN; presented only as commonly
reported.\cite{s-legacy}
\end{itemize}

\subsection{CJK-market office suites: Ichitaro and Hancom}\label{sec:cjkoffice}

\begin{itemize}[leftmargin=1.2em,itemsep=2pt]
\item \textbf{JustSystems Ichitaro 一太郎 (Japan; first released 1985; the \#2 Japanese
word processor after Word).} Default Japanese body font was \textbf{MS Mincho through
Ichitaro 2020}; per JustSystems' support note \#058017 and contemporaneous trade-press
reporting, \textbf{Ichitaro 2021 (released 2021-02-05) switched the default to Yu Mincho
\ja{游明朝}}---mirroring Office 2016's JP switch (\S\ref{sec:msoffice}) five years later
(the vendor note itself was not fetched during this research; attribution rests on the
reporting). Premium editions bundle licensed professional fonts (Morisawa faces,
Chikushi 筑紫 families, UD fonts).\cite{s-ichitaro,s-ichitaro-2,s-ichitaro-3}
\item \textbf{Hancom Office / Hangul HWP (Korea; dominant in Korean government).} The
historical Korean default serif category is \textbf{\ko{바탕} Batang} (the 1993 Ministry
of Culture standardization renamed the category from \emph{myeongjo}); recent Hancom
versions default to Hancom's own \textbf{\ko{함초롬바탕} HCR Batang} (Hamchorom family;
only HCR Batang/Dotum are licensed for use outside Hancom Office). The exact HWP release
that switched to Hamchorom is UNCERTAIN (Hangul 2010/2014 era).\cite{s-hancom,s-hancom-2}
\end{itemize}

% === §5 PDF and scan/OCR tools ===============================================
\section{PDF and Scan/OCR Tools}\label{sec:pdf}

\subsection{Adobe Acrobat and Reader}\label{sec:acrobat}

Figure~\ref{fig:acrobat} traces the Adobe Acrobat font-substitution provenance markers.

\begin{figure*}[t]
\centering
\begin{fptimeline}
\fpevent{1993}{Acrobat 1.0 and PDF 1.0 (Jun 15)---Base-14 standard fonts guaranteed without embedding}
\fpevent{1994}{Acrobat 2.0---Adobe Sans MM and Adobe Serif MM multiple-master substitution fonts introduced}
\fpevent{2006}{Acrobat 8 era---Typewriter tool defaults to Courier; form fields Helvetica 12}
\fpevent{2011}{Acrobat X and later---Edit PDF and TouchUp fall back to Minion Pro for missing Latin fonts}
\end{fptimeline}
\caption{Adobe Acrobat font-substitution provenance markers.}
\label{fig:acrobat}
\end{figure*}

\begin{itemize}[leftmargin=1.2em,itemsep=2pt]
\item \textbf{1993-06-15 --- the Base-14.} PDF 1.0 (announced at Comdex Fall 1992;
Acrobat 1.0 shipped June 1993) guarantees \textbf{14 standard Type 1 fonts} every
conforming viewer supplies without embedding: Times, Helvetica, and Courier in four
styles each, plus Symbol and ZapfDingbats (normative in the PDF Reference and ISO
32000-1 \S9.6.2.2). Viewers map them to system faces (Helvetica$\rightarrow$Arial,
Times$\rightarrow$Times New Roman, Courier$\rightarrow$Courier New on Windows).
\emph{Provenance:} unembedded Base-14 face names mark programmatically-generated or
default-styled PDFs relying on viewer
substitution.\cite{p-base14,p-base14-2,p-base14-3,p-history,p-history-2}
\item \textbf{1994 --- Acrobat 2.0: Adobe Sans MM / Adobe Serif MM.} Two multiple-master
fonts distributed only inside Acrobat/Reader (not installable or licensable separately),
used to synthesize metric-matched stand-ins for missing non-Base-14 Latin fonts.
\emph{Provenance:} their appearance in rendered or re-saved output is strongly consistent
with Acrobat processing a PDF whose fonts were absent---the strength of the inference
comes from their Acrobat-only distribution, though other Adobe-derived workflows can in
principle carry the names.\cite{p-mm,p-mm-2}
\item \textbf{Tool defaults.} Typewriter tool: \textbf{Courier} (Acrobat 8 era), later
Helvetica; form text fields: \textbf{Helvetica 12}; Fill \& Sign: fixed by design. Edit
PDF/TouchUp fallback for unavailable Latin fonts: \textbf{Minion Pro} (Acrobat X/XI/DC;
configurable via Content Editing preferences). \emph{Provenance:} Minion Pro body text
inside an otherwise foreign document marks post-hoc Acrobat
editing.\cite{p-tools,p-tools-2,p-minion,p-minion-2}
\item \textbf{CJK substitution packs.} Acrobat's Asian font packs supply \textbf{Kozuka
Mincho/Gothic} (JP), \textbf{Adobe Song Std} (SC), \textbf{Adobe Ming Std} (TC),
\textbf{Adobe Myungjo Std} (KR) as substitution faces for unembedded CJK fonts (pack
composition evolved across Acrobat 7$\rightarrow$DC; dates approximate). \emph{Provenance:}
CJK text rendering in these Adobe Std/Kozuka faces fingerprints Acrobat substitution
rather than original typesetting.\cite{p-cjk,p-cjk-2,p-cjk-3}
\end{itemize}

\subsection{Scan/OCR tools (PaperPort and peers)}\label{sec:ocr}

Provenance principle: \textbf{an OCR'd document's text layer carries the OCR engine's
output-font choices, not the original's typography.}

\begin{itemize}[leftmargin=1.2em,itemsep=2pt]
\item \textbf{ABBYY FineReader} (best documented): assigns recognized text to
visually-closest fonts by default and supports forcing a single output font---commonly
Arial or Times New Roman; ABBYY's Linux engine documentation expects Arial/Times New
Roman present. A uniform Arial/Times text layer over a scan is consistent with
ABBYY-class OCR output.\cite{p-abbyy,p-abbyy-2,p-abbyy-3}
\item \textbf{PaperPort} (Visioneer 1992 $\rightarrow$ ScanSoft 1999 $\rightarrow$ Nuance
2005 $\rightarrow$ Kofax 2018 $\rightarrow$ Tungsten Automation 2024): distributed at
scale as SOHO scanner bundleware---licensed with Brother, Canon, Dell, HP, and Xerox
devices per the bundling model described in Nuance's SEC filings (which evidence the
\emph{distribution channel}; no installed-base or market-share figures were located, so
``popularity'' is not quantified here). Historically carried the OmniPage OCR engine. Its
OCR output-font behavior is \textbf{undocumented}---stated as such, not
invented.\cite{p-paperport,p-paperport-2,p-paperport-3}
\item \textbf{OmniPage} (Caere, late 1980s $\rightarrow$ ScanSoft 2000 $\rightarrow$
Nuance $\rightarrow$ Tungsten) and \textbf{Readiris} (I.R.I.S., Belgium, 1987+, now a
Canon company): same class; output-font defaults undocumented---UNCERTAIN.\cite{p-omnipage,p-omnipage-2}
\item \textbf{Other PDF editors:} Foxit's add-text defaults to Helvetica; PDF-XChange's
Typewriter historically Courier New 12; Nitro exposes no documented single default
(UNCERTAIN).\cite{p-editors,p-editors-2,p-editors-3}
\end{itemize}

% === §6 Document-generation libraries ========================================
\section{Document-Generation Libraries (Aspose and Peers)}\label{sec:libraries}

Documents produced or converted by server-side libraries carry the library's default
and substitution fonts---a distinctive toolchain fingerprint.

\begin{table*}[t]
\centering\small
\caption{Default and missing-font behavior of the Aspose document-generation suites.}
\label{tab:aspose}
\renewcommand{\arraystretch}{1.25}
\begin{tabular}{@{}>{\raggedright\arraybackslash}p{0.13\textwidth} >{\raggedright\arraybackslash}p{0.17\textwidth} >{\raggedright\arraybackslash}p{0.30\textwidth} >{\raggedright\arraybackslash}p{0.30\textwidth}@{}}
\toprule
Product & New-object default & Missing-font behavior & Provenance tell \\
\midrule
Aspose.Words & Times New Roman 12 (Normal style) & Substitution chain ends at \code{DefaultFontName} = ``Times New Roman''; then first available font; last resort: embedded \textbf{Fanwood.ttf} & Fanwood output, with Aspose's substitution warning or matching metadata, is strong evidence of Aspose.Words on a font-starved host (Fanwood alone is only consistent with it---the font is publicly available) \\
\addlinespace
Aspose.Cells & Calibri 11 (workbook default style) & Falls back to workbook \code{DefaultStyle.Font}; PDF/image save options can override & Calibri mirrors Excel 2007+ \\
\addlinespace
Aspose.Slides & none fixed (deck theme fonts) & Heuristic closest-available system font, PowerPoint-like; ships no fonts & No single fingerprint; Cambria Math irreplaceable for equations \\
\addlinespace
Aspose.PDF & Times New Roman (TextFragment default) & Exact-name match, then \code{FontRepository.Substitutions}; Base-14 never embedded & Times New Roman re-set on edited text \\
\bottomrule
\end{tabular}
\end{table*}

Key documented facts: Aspose.Words' \code{DefaultFontSubstitutionRule.\allowbreak{}DefaultFontName}
documents ``The default value is `Times New Roman'\,''; the substitution order (exact
match $\rightarrow$ embedded $\rightarrow$ AltName $\rightarrow$ rules $\rightarrow$ first
available $\rightarrow$ bundled Fanwood) is Aspose's own documentation, including the
characteristic warning ``Font `Times New Roman' has not been found. Using `Fanwood' font
instead.''; Aspose.Slides' documentation states Aspose distributes no fonts and picks the
closest system font. None of the products ship CJK fonts---CJK output on bare containers
produces tofu or host-dependent fallback (the ordering is \emph{not} Aspose-defined), and
Aspose.PDF throws on standard CJK PostScript names (STSong-Light, MSung-Light,
HYSMyeongJo-Medium) when absent. \textbf{Stability finding (scoped):} in the Aspose
release notes and documentation searched (July 2026, keyed on default-font changes), no
change to these defaults was found---in particular nothing tracking Microsoft's 2024 Aptos
switch. Within that search scope the defaults appear stable, which cuts both ways: an
Aspose default alone cannot narrow a document to a release window, and an Aptos-defaulted
document is inconsistent with a factory-default Aspose
path.\cite{as-words,as-words-2,as-words-3,as-cells,as-slides,as-pdf}

Peers, for contrast: \textbf{iText} defaults to Helvetica (Base-14 metrics bundled, never
embedded); \textbf{Apache PDFBox} has no implicit default (caller chooses);
\textbf{python-docx} inherits Calibri 11 from its bundled Word template; \textbf{openpyxl}
hardcodes Calibri 11 (\code{DEFAULT\_FONT} in \code{styles/fonts.py}); \textbf{GemBox}
products default to Calibri.\cite{as-peers}

% === §7 Cross-platform provenance anchors ====================================
\section{Cross-Platform Provenance Anchors (Quick Reference)}\label{sec:anchors}

Selected anchors, each typed---because ``first design'', ``first public beta'', ``first
OS bundling'', and ``first default binding'' are different events with different
evidential force. Anachronism reasoning must use the \textbf{earliest public availability
of any kind, including public betas and standalone font releases}; the later anchor types
(bundling, default binding) support consistency reasoning about producing toolchains, not
impossibility claims. See the per-section registers for citations (Table~\ref{tab:anchors}).

\begin{table*}[t]
\centering\footnotesize
\caption{Cross-platform provenance anchors, each typed by anchor kind.}
\label{tab:anchors}
\renewcommand{\arraystretch}{1.2}
\begin{tabular}{@{}>{\raggedright\arraybackslash}p{0.235\textwidth} >{\raggedright\arraybackslash}p{0.15\textwidth} >{\raggedright\arraybackslash}p{0.26\textwidth} >{\raggedright\arraybackslash}p{0.20\textwidth}@{}}
\toprule
Font & Date & Anchor type & Vehicle \\
\midrule
Arial, Times New Roman, Courier New (TrueType) & 1992-04-06 & first OS bundling & Windows 3.1 \\
Comic Sans MS & 1995 & first bundling (add-on pack/OEM) & Windows 95 Plus! \\
Verdana, Georgia & 1996 & font release (IE/web-fonts program); OS-bundled by 98/2000 & Microsoft web fonts \\
Tahoma & 1994--1995 (Office/IE); 2000-02-17 & font release; then first OS default binding & Windows 2000 \\
Lucida Grande & 2000-09-13 & first OS bundling and default (public beta) & Mac OS X Public Beta \\
Hiragino (on Mac) & 2001-03-24 & first Mac OS bundling (existed earlier as a Jiyukobo/SCREEN commercial family) & Mac OS X 10.0 \\
Segoe UI & 2006-11-30 / 2007-01-30 & first OS bundling and default & Windows Vista \\
Calibri and the ClearType six & 2006-05-23; GA 2007-01-30 & first public availability (Office 2007 Beta 2); then GA & Office 2007 / Vista \\
Microsoft YaHei, JhengHei, Meiryo, Malgun Gothic & 2007-01-30 & first OS bundling and per-locale default & Windows Vista \\
Droid Sans & 2008-09-23 (announced 2007-11-12) & first OS bundling and default & Android 1.0 \\
Ubuntu Font Family & 2010-10-10 & first OS bundling and default & Ubuntu 10.10 \\
Cantarell & 2011-04-06 & first desktop default binding (font designed 2009) & GNOME 3.0 \\
Roboto & 2011-10-19 (SDK) / Nov 2011 (device) & first OS bundling and default & Android 4.0 \\
Yu Gothic / Yu Mincho & 2013-10-17 & first OS bundling (existed earlier as Jiyukobo commercial faces) & Windows 8.1 \\
Source Han Sans / Noto Sans CJK & 2014-07-15 & font release (open source) & Adobe/Google \\
San Francisco (SF) & 2014-11-18 (developer) / 2015-04-24 (watch) / Sep 2015 (iOS 9, El Capitan) & developer release; then OS default & Apple platforms \\
PingFang SC/TC/HK & 2015-09-16 (iOS 9) / 2015-09-30 (El Capitan) & first OS bundling and default & Apple platforms \\
DengXian 等线 & 2015-07-29 & first OS availability (Windows 10 FOD); Office 2016 SC default binding later & Microsoft \\
Segoe UI Variable, Cascadia Code & 2021-10-05 & first OS bundling (Cascadia Code open-sourced 2019) & Windows 11 \\
Roboto Flex & 2022-05-05 & font release (never an AOSP default) & Google Fonts \\
Aptos & 2023-07-13 & app default announcement, then staged rollout (M365 cloud fonts) & Microsoft 365 \\
Adwaita Sans & 2025-03-19 & first desktop default binding (Inter, its basis, released 2016) & GNOME 48 \\
\bottomrule
\end{tabular}
\end{table*}

% === §8 Limitations ==========================================================
\section{Limitations}\label{sec:limitations}

\begin{enumerate}[leftmargin=1.5em,itemsep=3pt]
\item \textbf{Verification asymmetry.} Every citation carries its verification tier;
[secondary] and UNCERTAIN labels are load-bearing. The residual weak points are
enumerated in place: pre-2000 CJK regional-edition ship dates (Microsoft's records begin
at Windows 2000); classic-Mac-era dates; the SF Pro rebrand year (2017, pending a Wayback
body read); several Wayback snapshot bodies confirmed to exist but unread due to
archive.org rate limiting during the research window; the explicitly unverified legacy
office defaults (\S\ref{sec:wps}); and the Aspose.Cells/Aspose.PDF and
iText/python-docx/GemBox default claims, whose specific source pages were not captured and
which are marked UNCERTAIN in their footnotes.
\item \textbf{Default-font attribution has no base rates.} No empirical evidence is
offered here for how often users, organizations, templates, or converters retain factory
defaults. Default-font matches therefore support only \emph{consistent-with} statements
about candidate toolchains---never probabilistic attribution, and phrases elsewhere in
this paper are calibrated accordingly. Anchor dates give \emph{terminus post quem} logic
only in the direction ``font used before its earliest public availability $\Rightarrow$
anachronism''; the converse direction is not symmetric.
\item \textbf{Availability precedes release for insiders.} Public-release anchors bound
\emph{ordinary} acquisition. Designers, foundries, beta testers, OEMs, and licensees hold
fonts earlier; where public betas materially widen early access (Calibri), the beta date
is the operative anchor, and \S\ref{sec:anchors} types each anchor for this reason.
\item \textbf{Presence $\neq$ bundling} (modern OSes). Windows 10+ Feature-on-Demand,
macOS download-on-demand, iOS 13+ user font installation, and Android's Downloadable
Fonts all decouple a font's presence on a device from the OS release that ``shipped'' it.
\item \textbf{Single-artifact inspections generalize narrowly.} The Office-theme
inspection (\S\ref{sec:msoffice}) is documented with path, version, date, and SHA-256, and
is valid for that build and channel only.
\item \textbf{Scope.} Timelines cover public/general releases; insider builds, betas, and
enterprise LTSC variants are noted only where they matter (Office 2007 Beta 2, iOS 7 beta
3). ``Current'' statements are bounded by the July 2026 research window.
\end{enumerate}

% =============================================================================
% unsrt (plain numeric) rather than unsrtnat: the paper cites numerically
% (natbib [numbers]), and unsrtnat's a--z author-year disambiguation overflows
% past 'z' with >26 same-author/year vendor-doc entries (e.g. Microsoft Learn),
% emitting an unbalanced brace into the .bbl. unsrt disambiguates by number, so
% [n] citations render identically and the .bbl stays well-formed.
\bibliographystyle{unsrt}
\bibliography{font-provenance}

\end{document}
